\section{Gaussian integrals and oscillator identities}\label{app:C.6}

When dark matter scatters off a nucleus in a crystal, the target operator is $\ee^{\ii\vect{q}\cdot\vect{u}}$, the exponential of a displacement that is a sum of harmonic-oscillator coordinates. Its thermal average gives the Debye--Waller factor, its matrix elements give the probabilities of creating zero, one or many phonons, and both reduce to Gaussian integrals or to their operator analogues. This section collects the identities that \cref{sec:5.3,sec:5.6,sec:6.1} use: Gaussian integrals in one and many variables, the Baker--Campbell--Hausdorff formula in the form needed for linear operators, the displacement operator and coherent states, and the statement $\langle\ee^{\op{A}}\rangle = \ee^{\langle \op{A}^2\rangle/2}$ for operators $\op{A}$ linear in the oscillator variables.

\subsection{Gaussian integrals}

For $\Real\mu > 0$ and any complex $\eta$,
\begin{equation}
  \int_{-\infty}^{\infty}\dd x\;\ee^{-\mu x^2 + \eta x} = \sqrt{\frac{\pi}{\mu}}\;\ee^{\eta^2/4\mu} .
  \label{eq:C.6-gauss}
\end{equation}
For $\eta = 0$ and real $\mu$, squaring the integral and changing to polar coordinates in the plane gives $\int\dd x\,\dd y\,\ee^{-\mu(x^2 + y^2)} = 2\pi\int_0^\infty\dd r\,r\,\ee^{-\mu r^2} = \pi/\mu$. For $\eta \neq 0$, completing the square, $-\mu x^2 + \eta x = -\mu(x - \eta/2\mu)^2 + \eta^2/4\mu$, and shifting the variable by $\eta/2\mu$ gives the result. For complex $\eta$ the shift moves the path of integration off the real axis; the integrand is analytic and decays at the ends of the strip in between, so by Cauchy's theorem the path can be moved back. The same argument extends the result to complex $\mu$ with $\Real\mu > 0$, with the branch of $\sqrt{\mu}$ that is positive for real $\mu$. The limit $\Real\mu \to 0^+$ gives the Fresnel integral, $\int\dd x\,\ee^{\ii\mu x^2} = \sqrt{\pi/\mu}\,\ee^{\ii\pi/4}$ for real $\mu > 0$, which appears in the time integrals of a freely recoiling nucleus (\cref{sec:11.4}).

With $\eta = \ii q$ and $\mu = 1/(2s^2)$, the Fourier transform of a normalized Gaussian of width $s$ is $\ee^{-q^2s^2/2}$, in any number of dimensions, as for the skin factor of the Helm form factor (\cref{sec:2.2}). Differentiating \cref{eq:C.6-gauss} at $\eta = 0$ with respect to $\mu$ gives the moments of the normalized distribution $\sqrt{\mu/\pi}\,\ee^{-\mu x^2}$,
\begin{equation}
  \langle x^{2n}\rangle = (2n - 1)!!\,\langle x^2\rangle^n, \qquad \langle x^2\rangle = \frac{1}{2\mu} ,
  \label{eq:C.6-moments}
\end{equation}
with odd moments zero. Reading \cref{eq:C.6-gauss} divided by $\sqrt{\pi/\mu}$ as an average gives the generating function
\begin{equation}
  \langle\ee^{\eta x}\rangle = \ee^{\eta^2\langle x^2\rangle/2} ,
  \label{eq:C.6-generating}
\end{equation}
which is the classical form of the operator identity of \cref{eq:C.6-bloch} below. Expanding both sides in powers of $\eta$ reproduces \cref{eq:C.6-moments}.

Let $\mathbf{M}$ be a real symmetric positive-definite $n\times n$ matrix and $\vect{s}$ a vector. In $n$ variables,
\begin{equation}
  \int\dd^nx\;\ee^{-\frac12\vect{x}\cdot\mathbf{M}\cdot\vect{x} + \vect{s}\cdot\vect{x}} = \frac{(2\pi)^{n/2}}{(\det\mathbf{M})^{1/2}}\;\ee^{\frac12\vect{s}\cdot\mathbf{M}^{-1}\cdot\vect{s}} .
  \label{eq:C.6-gaussn}
\end{equation}
The proof rotates to the eigenvectors of $\mathbf{M}$, which leaves $\dd^nx$ unchanged and turns the integral into a product of $n$ copies of \cref{eq:C.6-gauss} with $\mu$ equal to half an eigenvalue. The product of the eigenvalues is $\det\mathbf{M}$, and the exponents combine into $\tfrac12\vect{s}\cdot\mathbf{M}^{-1}\cdot\vect{s}$. For the normalized distribution, differentiating twice with respect to $\vect{s}$ at $\vect{s} = 0$ gives the covariance $\langle x_ix_j\rangle = (\mathbf{M}^{-1})_{ij}$, and differentiating four times gives
\begin{equation}
  \langle x_ix_jx_kx_l\rangle = \langle x_ix_j\rangle\langle x_kx_l\rangle + \langle x_ix_k\rangle\langle x_jx_l\rangle + \langle x_ix_l\rangle\langle x_jx_k\rangle ,
  \label{eq:C.6-isserlis}
\end{equation}
the sum over the three ways of pairing the four variables~\cite{isserlis1918formula}. Every higher even moment is the sum over all pairings in the same way, because the generating function is the exponential of a quadratic form in $\vect{s}$. In quantum many-body theory this statement is Wick's theorem (\cref{sec:4.8}).

\subsection{Operator identities}

Products of exponentials of operators need care, because $\ee^{\op{X}}\ee^{\op{Y}} \neq \ee^{\op{X} + \op{Y}}$ when $\op{X}$ and $\op{Y}$ do not commute. We first need the transformation of an operator $\op{Y}$ by $\ee^{\op{X}}$. The function $f(\lambda) = \ee^{\lambda \op{X}}\op{Y}\ee^{-\lambda \op{X}}$ obeys $f'(\lambda) = \ee^{\lambda \op{X}}[\op{X}, \op{Y}]\ee^{-\lambda \op{X}}$, and repeating the step gives the Taylor series
\begin{equation}
  \ee^{\op{X}}\op{Y}\ee^{-\op{X}} = \op{Y} + [\op{X}, \op{Y}] + \frac{1}{2!}[\op{X}, [\op{X}, \op{Y}]] + \frac{1}{3!}[\op{X}, [\op{X}, [\op{X}, \op{Y}]]] + \cdots .
  \label{eq:C.6-hadamard}
\end{equation}
We also need the Baker--Campbell--Hausdorff formula for the case in which the commutator $[\op{X}, \op{Y}]$ commutes with both $\op{X}$ and $\op{Y}$, as it does when $\op{X}$ and $\op{Y}$ are linear in boson operators and $[\op{X}, \op{Y}]$ is a number. Let $g(\lambda) = \ee^{\lambda \op{X}}\ee^{\lambda \op{Y}}$. Then $g'(\lambda) = (\op{X} + \ee^{\lambda \op{X}}\op{Y}\ee^{-\lambda \op{X}})\,g(\lambda) = (\op{X} + \op{Y} + \lambda[\op{X}, \op{Y}])\,g(\lambda)$, where the series \cref{eq:C.6-hadamard} stops after the commutator. The operators $\op{X} + \op{Y}$ and $[\op{X}, \op{Y}]$ commute, so the solution with $g(0) = 1$ is $g(\lambda) = \exp[\lambda(\op{X} + \op{Y}) + \tfrac12\lambda^2[\op{X}, \op{Y}]]$, and at $\lambda = 1$
\begin{equation}
  \ee^{\op{X}}\ee^{\op{Y}} = \ee^{\op{X} + \op{Y}}\,\ee^{[\op{X}, \op{Y}]/2}, \qquad \ee^{\op{X}}\ee^{\op{Y}} = \ee^{\op{Y}}\ee^{\op{X}}\,\ee^{[\op{X}, \op{Y}]} .
  \label{eq:C.6-BCH}
\end{equation}
The second form follows from the first by exchanging $\op{X}$ and $\op{Y}$ and dividing. For general $\op{X}$ and $\op{Y}$ the exponent contains nested commutators to all orders~\cite{wilcox1967exponential}; they are not needed here.

\subsection{Displacement operators and coherent states}

We write the oscillator as in \cref{we:2-oscillator}, $H = \omega_0(d^\dagger d + \tfrac12)$ with $[d, d^\dagger] = 1$, and its coordinate as $x = \ell_0(d + d^\dagger)$ with the zero-point length $\ell_0 = (2M\omega_0)^{-1/2}$ for a mass $M$; for a phonon mode it is $\ell_{\vect{Q}\nu}$ of \cref{eq:A.6-ell}. The displacement operator with a complex amplitude $\alpha_{\mathrm{d}}$ is
\begin{equation}
  \op{D}(\alpha_{\mathrm{d}}) = \exp\bigl(\alpha_{\mathrm{d}}d^\dagger - \alpha_{\mathrm{d}}^*d\bigr) = \ee^{-|\alpha_{\mathrm{d}}|^2/2}\,\ee^{\alpha_{\mathrm{d}}d^\dagger}\ee^{-\alpha_{\mathrm{d}}^*d} ,
  \label{eq:C.6-D}
\end{equation}
where the second form is \cref{eq:C.6-BCH} with $\op{X} = \alpha_{\mathrm{d}}d^\dagger$, $\op{Y} = -\alpha_{\mathrm{d}}^*d$ and $[\op{X}, \op{Y}] = |\alpha_{\mathrm{d}}|^2$. It is unitary, and by \cref{eq:C.6-hadamard} it shifts the ladder operators, $\op{D}^\dagger d\,\op{D} = d + \alpha_{\mathrm{d}}$, so it shifts the coordinate by $2\ell_0\Real\alpha_{\mathrm{d}}$ and the momentum in proportion to $\Imag\alpha_{\mathrm{d}}$. A momentum kick is the special case of a real coordinate: $\ee^{\ii qx} = \ee^{\ii q\ell_0(d + d^\dagger)} = \op{D}(\ii q\ell_0)$, which displaces the momentum by $q$. This is the operator through which dark matter acts on a bound nucleus.

Acting on the ground state, $\ee^{-\alpha_{\mathrm{d}}^*d}|0\rangle = |0\rangle$, and $\ee^{\alpha_{\mathrm{d}}d^\dagger}|0\rangle = \sum_n\alpha_{\mathrm{d}}^n(d^\dagger)^n|0\rangle/n!$, which is $\sum_n\alpha_{\mathrm{d}}^n|n\rangle/\sqrt{n!}$ by $(d^\dagger)^n|0\rangle = \sqrt{n!}\,|n\rangle$ (\cref{sec:2.3}). The result is a coherent state,
\begin{equation}
  \op{D}(\alpha_{\mathrm{d}})|0\rangle = \ee^{-|\alpha_{\mathrm{d}}|^2/2}\sum_{n=0}^{\infty}\frac{\alpha_{\mathrm{d}}^n}{\sqrt{n!}}\,|n\rangle, \qquad
  |\langle n|\op{D}(\alpha_{\mathrm{d}})|0\rangle|^2 = \ee^{-|\alpha_{\mathrm{d}}|^2}\frac{|\alpha_{\mathrm{d}}|^{2n}}{n!} .
  \label{eq:C.6-coherent}
\end{equation}
The number of quanta in a displaced ground state is Poisson distributed with mean $|\alpha_{\mathrm{d}}|^2$. For a kick, $|\alpha_{\mathrm{d}}|^2 = q^2\ell_0^2 = q^2/(2M\omega_0)$, the recoil energy $q^2/2M$ in units of the quantum $\omega_0$. The probability of leaving the oscillator in its ground state is $\ee^{-q^2\ell_0^2}$, and the mean energy deposited is $\omega_0|\alpha_{\mathrm{d}}|^2 = q^2/2M$, the free recoil energy, at every $q$. For a \SI{60}{meV} mode of a silicon nucleus, $M = \SI{26.16}{GeV}$ (the mean atomic mass of natural silicon, as in \cref{sec:5.3}), $q^2\ell_0^2 = 1$ at $q = (2M\omega_0)^{1/2} = \SI{56}{keV}$: below it a kick mostly creates no quantum or one, and far above it the Poisson distribution is broad and the nucleus recoils almost as if free. The same distribution, with $|\alpha_{\mathrm{d}}|^2$ the Huang--Rhys factor $S_{\mathrm{HR}}$, describes an electronic transition that displaces the equilibrium position of an oscillator (\cref{sec:5.6}).

Ordering an exponential so that all annihilation operators stand to the right makes thermal and ground-state averages easy. For a real or complex $\eta$, \cref{eq:C.6-BCH} with $\op{X} = \eta d^\dagger$, $\op{Y} = \eta d$ and $[\op{X}, \op{Y}] = -\eta^2$ gives
\begin{equation}
  \ee^{\eta(d + d^\dagger)} = \ee^{\eta d^\dagger}\ee^{\eta d}\,\ee^{\eta^2/2} .
  \label{eq:C.6-normal}
\end{equation}
In the ground state, $\langle0|\ee^{\eta d^\dagger}\ee^{\eta d}|0\rangle = 1$, so $\langle0|\ee^{\eta(d + d^\dagger)}|0\rangle = \ee^{\eta^2/2}$. For a kick, $\eta = \ii q\ell_0$ and the ground-state average is $\ee^{-q^2\ell_0^2/2}$, the square root of the ground-state survival probability, as it must be. Since $\langle0|(d + d^\dagger)^2|0\rangle = 1$, the exponent is $\tfrac12\langle0|(\eta(d + d^\dagger))^2|0\rangle$: the ground-state average of the exponential is the exponential of half the average of the square, as for the classical Gaussian of \cref{eq:C.6-generating}.

\subsection{Thermal averages of exponentials}

The ground state of an oscillator is a Gaussian in $x$, and so is its thermal state: the density matrix $\ee^{-\beta H}/\mathcal{Z}$ is Gaussian in the coordinate and the momentum. The identity of the last paragraph therefore extends to any temperature and to any number of independent oscillators. For an operator $\op{A}$ that is a linear combination of the $d_\gamma$ and $d^\dagger_\gamma$ of a set of harmonic modes in thermal equilibrium,
\begin{equation}
  \langle\ee^{\op{A}}\rangle = \ee^{\langle \op{A}^2\rangle/2} ,
  \label{eq:C.6-bloch}
\end{equation}
which is Bloch's identity, proved in \cref{app:D.15}; Mermin gave a short alternative derivation~\cite{mermin1966short}. In the classical limit $T \gg \omega_0$ it is \cref{eq:C.6-generating}, since the coordinates are then classical Gaussian variables. For one mode and $\op{A} = \ii q x$, it gives $\langle\ee^{\ii qx}\rangle = \ee^{-q^2\langle x^2\rangle/2}$ with $\langle x^2\rangle = \ell_0^2(2n_{\mathrm{B}}(\omega_0) + 1)$, the zero-point value multiplied by the factor found in \cref{we:2-oscillator} and in the Matsubara sum of \cref{app:C.5}. For a crystal, $\op{A} = \ii\vect{q}\cdot\vect{u}_\kappa$ and the result is $\ee^{-W_\kappa(\vect{q})}$ with $2W_\kappa(\vect{q}) = \vect{q}\cdot B_\kappa\cdot\vect{q}$, the Debye--Waller factor of \cref{eq:A.6-DW} at the level of the amplitude.

Combined with \cref{eq:C.6-BCH}, Bloch's identity gives the average of a product. For $\op{A}$ and $\op{B}$ both linear, $[\op{A}, \op{B}]$ is a number, $\ee^{\op{A}}\ee^{\op{B}} = \ee^{\op{A} + \op{B}}\ee^{[\op{A}, \op{B}]/2}$, and $\langle\ee^{\op{A} + \op{B}}\rangle = \ee^{\langle(\op{A} + \op{B})^2\rangle/2}$. Expanding $(\op{A} + \op{B})^2 = \op{A}^2 + \op{B}^2 + \op{A} \op{B} + \op{B} \op{A}$ and adding $\tfrac12[\op{A}, \op{B}] = \tfrac12\langle \op{A} \op{B} - \op{B} \op{A}\rangle$ to the exponent,
\begin{equation}
  \langle\ee^{\op{A}}\ee^{\op{B}}\rangle = \exp\Bigl[\tfrac12\langle \op{A}^2\rangle + \tfrac12\langle \op{B}^2\rangle + \langle \op{A} \op{B}\rangle\Bigr] .
  \label{eq:C.6-product}
\end{equation}
With $\op{A} = -\ii\vect{q}\cdot\vect{u}_\kappa(t)$ and $\op{B} = \ii\vect{q}\cdot\vect{u}_\kappa(0)$, which are linear because the Heisenberg operators of a harmonic crystal are $d_\gamma(t) = d_\gamma\ee^{-\ii\omega_\gamma t}$, the first two terms give $\ee^{-2W_\kappa(\vect{q})}$ and the third gives $\ee^{\langle(\vect{q}\cdot\vect{u}_\kappa(t))(\vect{q}\cdot\vect{u}_\kappa(0))\rangle}$. This is the generating formula from which \cref{sec:5.6} builds the multiphonon response: expanding the last exponential in powers of the displacement correlation counts the phonons created, and the zeroth term is the elastic weight $\ee^{-2W_\kappa}$.
